% Options for packages loaded elsewhere
\PassOptionsToPackage{unicode}{hyperref}
\PassOptionsToPackage{hyphens}{url}
\documentclass[
  a4paper,11pt]{article}
\usepackage{xcolor}
\usepackage[margin=1in]{geometry}
\usepackage{amsmath,amssymb}
\setcounter{secnumdepth}{-\maxdimen} % remove section numbering
\usepackage{iftex}
\ifPDFTeX
  \usepackage[T1]{fontenc}
  \usepackage[utf8]{inputenc}
  \usepackage{textcomp} % provide euro and other symbols
\else % if luatex or xetex
  \usepackage{unicode-math} % this also loads fontspec
  \defaultfontfeatures{Scale=MatchLowercase}
  \defaultfontfeatures[\rmfamily]{Ligatures=TeX,Scale=1}
\fi
\usepackage{lmodern}
\ifPDFTeX\else
  % xetex/luatex font selection
\fi
% Use upquote if available, for straight quotes in verbatim environments
\IfFileExists{upquote.sty}{\usepackage{upquote}}{}
\IfFileExists{microtype.sty}{% use microtype if available
  \usepackage[]{microtype}
  \UseMicrotypeSet[protrusion]{basicmath} % disable protrusion for tt fonts
}{}
\makeatletter
\@ifundefined{KOMAClassName}{% if non-KOMA class
  \IfFileExists{parskip.sty}{%
    \usepackage{parskip}
  }{% else
    \setlength{\parindent}{0pt}
    \setlength{\parskip}{6pt plus 2pt minus 1pt}}
}{% if KOMA class
  \KOMAoptions{parskip=half}}
\makeatother
\usepackage{longtable,booktabs,array}
\newcounter{none} % for unnumbered tables
\usepackage{calc} % for calculating minipage widths
% Correct order of tables after \paragraph or \subparagraph
\usepackage{etoolbox}
\makeatletter
\patchcmd\longtable{\par}{\if@noskipsec\mbox{}\fi\par}{}{}
\makeatother
% Allow footnotes in longtable head/foot
\IfFileExists{footnotehyper.sty}{\usepackage{footnotehyper}}{\usepackage{footnote}}
\makesavenoteenv{longtable}
\setlength{\emergencystretch}{3em} % prevent overfull lines
\providecommand{\tightlist}{%
  \setlength{\itemsep}{0pt}\setlength{\parskip}{0pt}}
\usepackage{bookmark}
\IfFileExists{xurl.sty}{\usepackage{xurl}}{} % add URL line breaks if available
\urlstyle{same}
\hypersetup{
  hidelinks,
  pdfcreator={LaTeX via pandoc}}

\author{}
\date{}

\begin{document}

\section{Geometric Necessity of Spontaneous Symmetry Breaking ---
Three-Stage Derivation of the Higgs
Mechanism}\label{geometric-necessity-of-spontaneous-symmetry-breaking-three-stage-derivation-of-the-higgs-mechanism}

\subsection{------ A Structural Answer to the Higgs Mechanism Problem of
{[}3{]} §9.11.7
------}\label{a-structural-answer-to-the-higgs-mechanism-problem-of-3-9.11.7}

\textbf{Author}: Noriaki Kihara (WF System Co., Ltd.)

\textbf{ORCID}: https://orcid.org/0009-0002-7981-4891

\textbf{DOI}: https://doi.org/10.5281/zenodo.19731606

\textbf{Date}: April 2026

\begin{center}\rule{0.5\linewidth}{0.5pt}\end{center}

\subsection{§0 Purpose of This Paper}\label{purpose-of-this-paper}

In {[}3{]} §9.11.7, it was declared that ``how the dynamical mechanism
of spontaneous symmetry breaking via the Higgs field is understood
within this framework remains as a future problem.''

In the Standard Model, this mechanism is explained by a Mexican-hat
potential \(V(\phi) = \mu^2 |\phi|^2 + \lambda |\phi|^4\) (with
\(\mu^2 < 0\)). The symmetric state \(\phi = 0\) is unstable, and
symmetry is spontaneously broken when the field falls into a non-zero
vacuum expectation value.

This paper demonstrates that within our framework, such an external
potential is unnecessary, and that spontaneous symmetry breaking has
already been derived as the following three-stage \textbf{chain of
geometric necessities}.

{\def\LTcaptype{none} % do not increment counter
\begin{longtable}[]{@{}
  >{\centering\arraybackslash}p{(\linewidth - 4\tabcolsep) * \real{0.3846}}
  >{\raggedright\arraybackslash}p{(\linewidth - 4\tabcolsep) * \real{0.3077}}
  >{\raggedright\arraybackslash}p{(\linewidth - 4\tabcolsep) * \real{0.3077}}@{}}
\toprule\noalign{}
\begin{minipage}[b]{\linewidth}\centering
Stage
\end{minipage} & \begin{minipage}[b]{\linewidth}\raggedright
Content
\end{minipage} & \begin{minipage}[b]{\linewidth}\raggedright
Source
\end{minipage} \\
\midrule\noalign{}
\endhead
\bottomrule\noalign{}
\endlastfoot
1 & Necessity of 4-dimensional spacetime & {[}2{]} \\
2 & Necessity of the 3+1 decomposition & {[}1{]} \\
3 & One-dimensional standing waves and the a posteriori nature of axis
labels & {[}4{]} \\
\end{longtable}
}

\textbf{Referenced papers}: - {[}1{]} Shape-Invariant Standing Waves on
Closed Spheres --- Stability by Dimension and the Geometric Necessity of
3+1 Spacetime - {[}2{]} Full Spherical Coverage of Central Projection in
Discrete Spaces and the Number-Theoretic Necessity of a 5-Dimensional
Background Space (W5), DOI: 10.5281/zenodo.19624957 - {[}3{]} Set
Structure and Combinatorial Properties of 6-Dimensional Hypercuboids
(W8), DOI: 10.5281/zenodo.19762134 - {[}4{]} Four Modes of
Shape-Invariant Waves (W10), DOI: 10.5281/zenodo.19709798

\begin{center}\rule{0.5\linewidth}{0.5pt}\end{center}

\subsection{§1 First Stage: Necessity of 4-Dimensional
Spacetime}\label{first-stage-necessity-of-4-dimensional-spacetime}

\subsubsection{1.1 Lagrange's Four-Square Theorem and Full Spherical
Coverage}\label{lagranges-four-square-theorem-and-full-spherical-coverage}

By {[}2{]}, the necessary and sufficient condition for full spherical
coverage of central projection on a discrete integer lattice
\(\mathbb{Z}^n\) is \(n = 5\) (the dimension of the background space).

This result depends on Lagrange's four-square theorem (every positive
integer can be expressed as the sum of four squares). For \(n = 4\),
full spherical coverage is not guaranteed, and for \(n \geq 6\), it is
redundant. The subjective space obtained by central projection from a
5-dimensional background space is 4-dimensional, and therefore:

\textbf{Proposition 1.1} ({[}2{]}). The minimal subjective space in
which full spherical coverage of central projection holds on a discrete
integer lattice is 4-dimensional.

\subsubsection{1.2 Consequence of This
Stage}\label{consequence-of-this-stage}

Four-dimensional spacetime is a geometric and number-theoretic
necessity, not an a priori assumption. Four coordinate axes
\(x, y, z, t\) exist. At this stage, all four axes are completely
equivalent, and the distinction between ``space'' and ``time'' does not
yet exist.

\begin{center}\rule{0.5\linewidth}{0.5pt}\end{center}

\subsection{§2 Second Stage: Necessity of the 3+1
Decomposition}\label{second-stage-necessity-of-the-31-decomposition}

\subsubsection{2.1 Stability Condition for Standing Waves on Closed
Spheres}\label{stability-condition-for-standing-waves-on-closed-spheres}

By {[}1{]}, the minimal dimension at which shape-invariant standing
waves can exist non-trivially and stably on a closed \(n\)-dimensional
sphere \(S^n\) is \(n = 3\). This result follows from the combination of
the following two conditions:

\begin{enumerate}
\def\labelenumi{(\roman{enumi})}
\item
  \textbf{Huygens' principle}: Wavefronts propagate without trailing
  tails, and the identity of wave packets is preserved, only on
  odd-dimensional spheres.
\item
  \textbf{Non-triviality}: \(S^1\) trivially satisfies the condition but
  does not exhibit non-trivial behavior involving wave-packet formation,
  restoration, and antipodal convergence.
\end{enumerate}

\textbf{Proposition 2.1} ({[}1{]}). The minimal dimension at which
shape-invariant standing waves can exist non-trivially and stably on a
closed \(n\)-dimensional sphere is \(n = 3\).

\subsubsection{2.2 Independent Argument via Frobenius'
Theorem}\label{independent-argument-via-frobenius-theorem}

By {[}1{]} §8, from Frobenius' theorem (the only finite-dimensional
associative division algebras over the real numbers are
\(\mathbb{R}, \mathbb{C}, \mathbb{H}\)), the maximum dimension in which
wave composition and interaction are algebraically closed is 4 (the
quaternions \(\mathbb{H}\)). The quaternions decompose intrinsically
within their algebraic structure into \(1 + 3\) (scalar part + vector
part).

\subsubsection{2.3 The 3+1 Decomposition of 4-Dimensional
Spacetime}\label{the-31-decomposition-of-4-dimensional-spacetime}

Combining the 4-dimensional spacetime of Proposition 1.1 with the
\(S^3\) stability condition of Proposition 2.1:

\textbf{Proposition 2.2}. For standing waves to exist stably in
4-dimensional spacetime, the four axes must decompose into 3-dimensional
space + 1-dimensional time.

\emph{Proof}. We exhaustively verify the decomposition schemes of 4
dimensions.

\begin{enumerate}
\def\labelenumi{(\roman{enumi})}
\item
  \textbf{4+0 (no decomposition)}: If all four axes remain equivalent,
  the closed space is \(S^4\) (the 4-dimensional sphere). Since \(S^4\)
  is even-dimensional, Huygens' principle does not hold, and standing
  waves cannot exist stably ({[}1{]} §6.3). Excluded.
\item
  \textbf{2+2 decomposition}: If 4 dimensions are decomposed into 2 + 2,
  the spatial closed submanifold is either \(S^2\) or
  \(S^1 \times S^1\). Since \(S^2\) is even-dimensional, Huygens'
  principle does not hold. \(S^1 \times S^1\), being a direct product of
  \(S^1\), does not possess non-trivial wave-packet structure ({[}1{]}
  §7.3). Neither satisfies the \(S^3\) stability condition of
  Proposition 2.1. Excluded.
\item
  \textbf{1+3 decomposition}: In the case of 1-dimensional space +
  3-dimensional time, the space becomes \(S^1\), which does not support
  non-trivial standing waves ({[}1{]} §7.3). Excluded.
\item
  \textbf{3+1 decomposition}: In the case of 3-dimensional space +
  1-dimensional time, the spatial closed submanifold is \(S^3\). Since
  \(S^3\) is odd-dimensional, Huygens' principle holds, and non-trivial,
  stable standing waves exist (Proposition 2.1). This is the
  \textbf{only admissible decomposition scheme}.
\end{enumerate}

Therefore, the decomposition scheme for stable existence of standing
waves in 4-dimensional spacetime is uniquely determined to be 3+1.
\(\square\)

\subsubsection{2.4 Consequence of This
Stage}\label{consequence-of-this-stage-1}

The 3+1 decomposition of 4-dimensional spacetime is a geometric
necessity derived from the stability condition for standing waves. At
this stage, the distinction between \(xyz\) (three spatial axes) and
\(t\) (one temporal axis) is established. The three spatial axes
\(x, y, z\) are completely equivalent under \(\mathrm{SO}(3)\) symmetry.

\begin{center}\rule{0.5\linewidth}{0.5pt}\end{center}

\subsection{§3 Third Stage: One-Dimensional Standing Waves and the A
Posteriori Nature of Axis
Labels}\label{third-stage-one-dimensional-standing-waves-and-the-a-posteriori-nature-of-axis-labels}

\subsubsection{3.1 The Meaning of κ=1}\label{the-meaning-of-ux3ba1}

By {[}4{]} §5, standing waves are modes with \(\kappa = 1\) (the wave
vector has a non-zero component along only one spatial axis). Here, it
is critically important to understand the precise meaning of
\(\kappa = 1\).

\(\kappa = 1\) refers to \textbf{an oscillation mode that has spatial
extent in only one dimension}. It does not mean that ``a specific axis
(say \(z\)) has been selected from among \(xyz\).'' In an
\(\mathrm{SO}(3)\)-symmetric 3-dimensional space, a one-dimensional
oscillation is physically equivalent regardless of which direction it
faces.

\subsubsection{3.2 The A Posteriori Nature of Axis
Labels}\label{the-a-posteriori-nature-of-axis-labels}

\textbf{Proposition 3.1}. When a \(\kappa = 1\) standing wave exists,
calling the spatial axis corresponding to the oscillation direction
``\(z\)'' is an a posteriori assignment of a coordinate label, not a
physical axis selection.

\emph{Proof}. By Proposition 5.1 of {[}4{]} §5.2, the waves with
\(\kappa = 1\) --- namely \(\mathbf{k} = (k_x, 0, 0, 0)\),
\(\mathbf{k} = (0, k_y, 0, 0)\), and \(\mathbf{k} = (0, 0, k_z, 0)\) ---
all exist equivalently by \(\mathrm{SO}(3)\) symmetry. These are
different coordinate representations of the same physical state, and
there is no dynamical problem of ``choosing one among three.''

After the standing wave exists, one simply names its oscillation
direction ``\(z\)'' a posteriori. \(\square\)

\subsubsection{3.3 Why One-Dimensional
Oscillation?}\label{why-one-dimensional-oscillation}

The fact that \(\kappa = 1\) is the unique scalar standing mode is
derived in Proposition 5.2 of {[}4{]} §5.4:

\begin{itemize}
\tightlist
\item
  \(\kappa = 0\): Spatially trivial (constant function) and possesses no
  phase structure
\item
  \(\kappa = 1\): Standing oscillation along one spatial axis. Scalar by
  \(\mathrm{SO}(3)\) averaging
\item
  \(\kappa = 2\): Intrinsically carries surface orientation (rotational
  structure) and is not a scalar
\item
  \(\kappa \geq 3\): Excluded by the stability condition of {[}4{]} §3
\end{itemize}

Therefore, spatially non-trivial scalar standing waves are restricted to
\(\kappa = 1\). Even if two- or three-dimensional oscillation modes were
to exist, their baseline (ground mode) reduces to a one-dimensional
oscillation.

In the Standard Model, the fact that the Higgs boson is a scalar
(spin-0) corresponds in our framework to \(\kappa = 1\) being the unique
scalar standing mode. That is, the correspondence ``Higgs = scalar =
\(\kappa = 1\) = one-dimensional oscillation'' is uniquely derived from
the structure of standing-wave modes.

\subsubsection{3.4 Consequence for Mass
Hierarchy}\label{consequence-for-mass-hierarchy}

The following mass hierarchy was already described in {[}3{]} §9.11.7,
but is summarized here to complete the three-stage derivation of this
supplementary lecture.

Once the oscillation direction of the standing wave has been called
\(z\):

\begin{itemize}
\tightlist
\item
  Fermions with a non-zero component along the \(z\) axis (a posteriori
  called the ``third generation'') directly share a spatial axis with
  the standing wave, resulting in maximal coupling → maximal mass
\item
  Fermions with non-zero components along the \(x\) and \(y\) axes (a
  posteriori called the ``first and second generations'') do not share a
  spatial axis with the standing wave, resulting in indirect coupling →
  lighter masses
\end{itemize}

This is precisely the generational mass hierarchy \(m_3 \gg m_{1,2}\)
described in {[}3{]} §9.11.7.

\begin{center}\rule{0.5\linewidth}{0.5pt}\end{center}

\subsection{§4 Integration of the Three Stages: Resolution of the
Problem}\label{integration-of-the-three-stages-resolution-of-the-problem}

\subsubsection{4.1 The Chain of Spontaneous Symmetry
Breaking}\label{the-chain-of-spontaneous-symmetry-breaking}

The dynamical mechanism of spontaneous symmetry breaking that {[}3{]}
§9.11.7 designated as a ``future problem'' has already been derived as
the following three stages of geometric necessity:

\textbf{First stage} ({[}2{]}): Full spherical coverage condition on
discrete integer lattice → \textbf{4-dimensional spacetime is necessary}

\[\text{$\mathbb{Z}^n$ full spherical coverage} \implies n_{\text{background}} = 5 \implies n_{\text{subjective}} = 4\]

\textbf{Second stage} ({[}1{]}): Stability condition for standing waves
→ \textbf{3+1 decomposition is necessary}

\[\text{Stable standing waves on $S^n$} \implies n = 3 \implies 4\text{D} = 3 + 1\]

\textbf{Third stage} ({[}4{]}): Uniqueness of scalar standing waves →
\textbf{\(\kappa = 1\) is necessary, and axis labels are a posteriori}

\[\text{Scalar standing wave} \implies \kappa = 1 \implies \text{the oscillation direction is called $z$}\]

\subsubsection{4.2 Dispensability of the Mexican-Hat
Potential}\label{dispensability-of-the-mexican-hat-potential}

\textbf{Proposition 4.1}. Within this framework, no external dynamical
mechanism such as a Mexican-hat potential is required for spontaneous
symmetry breaking.

\emph{Proof}. We verify in turn the three elements demanded by the Higgs
mechanism of the Standard Model:

\begin{enumerate}
\def\labelenumi{(\roman{enumi})}
\item
  \textbf{Existence of a symmetric initial state}: In the first stage,
  four axes exist on equal footing. In the second stage,
  \(\mathrm{SO}(3)\) symmetry among \(xyz\) is established. ✓
\item
  \textbf{Occurrence of symmetry breaking}: In the third stage, the
  \(\kappa = 1\) standing wave has a one-dimensional oscillation
  direction. This is not a dynamical selection but the very definition
  of the \(\kappa = 1\) mode. ✓
\item
  \textbf{Physical consequences after breaking (mass hierarchy)}: As
  derived in §3.4, the presence or absence of \(z\)-axis sharing
  determines the coupling strength, and \(m_3 \gg m_{1,2}\) follows. ✓
\end{enumerate}

All three elements follow from geometric necessity, requiring no
introduction of a potential. \(\square\)

\subsubsection{4.3 Comparison with the Standard
Model}\label{comparison-with-the-standard-model}

{\def\LTcaptype{none} % do not increment counter
\begin{longtable}[]{@{}
  >{\raggedright\arraybackslash}p{(\linewidth - 4\tabcolsep) * \real{0.3333}}
  >{\raggedright\arraybackslash}p{(\linewidth - 4\tabcolsep) * \real{0.3333}}
  >{\raggedright\arraybackslash}p{(\linewidth - 4\tabcolsep) * \real{0.3333}}@{}}
\toprule\noalign{}
\begin{minipage}[b]{\linewidth}\raggedright
Element
\end{minipage} & \begin{minipage}[b]{\linewidth}\raggedright
Standard Model
\end{minipage} & \begin{minipage}[b]{\linewidth}\raggedright
This framework
\end{minipage} \\
\midrule\noalign{}
\endhead
\bottomrule\noalign{}
\endlastfoot
Cause of symmetry breaking & \(\mu^2 < 0\) in \(V(\phi)\) & Definition
of \(\kappa = 1\) (one-axis oscillation) \\
Why the symmetric state is unstable & Shape of the potential & 4D → 3+1
decomposition is the stability condition ({[}1{]}) \\
Origin of the vacuum expectation value &
\(\langle\phi\rangle = v \neq 0\) & Ground mode of the standing wave
({[}4{]} §5.3) \\
Mechanism of mass acquisition & Yukawa coupling
\(y_f \langle\phi\rangle\) & Phase coupling on shared spatial axes
({[}3{]} §9.11.7) \\
Free parameters & \(\mu^2, \lambda\) & None (geometrically
determined) \\
\end{longtable}
}

\begin{center}\rule{0.5\linewidth}{0.5pt}\end{center}

\subsection{§5 Conclusion}\label{conclusion}

In this paper, we have provided the following structural answer to the
Higgs mechanism problem left unresolved in {[}3{]} §9.11.7.

\begin{enumerate}
\def\labelenumi{\arabic{enumi}.}
\item
  \textbf{Spontaneous symmetry breaking is a chain of three stages of
  geometric necessity.} Each stage --- the first stage (necessity of 4
  dimensions, {[}2{]}), the second stage (necessity of the 3+1
  decomposition, {[}1{]}), and the third stage (the a posteriori nature
  of axis labels for \(\kappa = 1\) standing waves, {[}4{]}) --- has
  already been derived in existing papers and requires no external
  dynamical mechanism.
\item
  \textbf{\(\kappa = 1\) does not mean ``selecting a specific axis'' but
  rather ``having spatial extent in only one dimension.''} In an
  \(\mathrm{SO}(3)\)-symmetric space, the direction of a one-dimensional
  oscillation is physically indistinguishable, and it is merely named
  ``\(z\)'' a posteriori.
\item
  \textbf{The Mexican-hat potential is dispensable.} The three elements
  that the Standard Model realizes through \(V(\phi)\) --- symmetry
  breaking, vacuum expectation value, and mass acquisition --- all
  follow from geometric necessity (Proposition 4.1).
\end{enumerate}

This supplementary lecture has derived the qualitative direction of the
mass hierarchy (\(m_3 \gg m_{1,2}\)) from geometric necessity, but the
quantitative derivation of specific mass ratios lies outside the scope
of this paper (see ``Analysis of Mass Structure --- A Framework for Mass
Analysis via Axis Scale Values and Signed Areas'').

\begin{center}\rule{0.5\linewidth}{0.5pt}\end{center}

\subsection{References}\label{references}

\begin{itemize}
\tightlist
\item
  {[}1{]} Noriaki Kihara, ``Shape-Invariant Standing Waves on Closed
  Spheres --- Stability by Dimension and the Geometric Necessity of 3+1
  Spacetime,'' 2026.
\item
  {[}2{]} Noriaki Kihara, ``Full Spherical Coverage of Central
  Projection in Discrete Spaces and the Number-Theoretic Necessity of a
  5-Dimensional Background Space,'' Zenodo preprint, DOI:
  10.5281/zenodo.19624957 (2026).
\item
  {[}3{]} Noriaki Kihara, ``Set Structure and Combinatorial Properties
  of 6-Dimensional Hypercuboids,'' Zenodo preprint, DOI:
  10.5281/zenodo.19762134 (2026).
\item
  {[}4{]} Noriaki Kihara, ``Four Modes of Shape-Invariant Waves ---
  Spin, Chirality, and Statistics Determined by Wave-Vector Structure,''
  Zenodo preprint, DOI: 10.5281/zenodo.19709798 (2026).
\end{itemize}

\end{document}
